\chapter{Introduction}
\label{ch:introduction}

\section{Background and Motivation}
undefined

\section{Research Problem and Objectives}
undefined

\section{Research Questions}
\begin{description}
    \item[\textbf{RQ1}] How reliably can Vericult assess arbitrary prompt--response pairs using the D01--D10 framework while distinguishing non-applicability, non-assessability and evidence insufficiency?
    \item[\textbf{RQ2}] In a controlled Best-of-4 setting, to what extent does Vericult identify the response preferred by human annotators, and how does this compare with an independent reward model and a direct LLM judge?
    \item[\textbf{RQ3}] How well does the frozen verifier transfer to independently published cultural-evaluation settings represented by the preregistered external protocol, including SafeWorld, CultureLLM/World Values Survey, ScopeBench and CARB where official item-level data are available?
    \item[\textbf{RQ4}] Which failure modes arise from applicability/assessability gating, context and dimension planning, target extraction, retrieval, evidence synthesis, selective fallback, abstention and final scoring, and what do these failures imply for future cultural-verification systems?
\end{description}

RQ1 reflects the intended single-response deployment scenario and explicitly treats the reject options as part of correct verifier behavior rather than as missing predictions. RQ2 uses Best-of-4 as a controlled preference-evaluation instrument and bridge to reward-model research. RQ3 tests transfer after the semantic freeze. RQ4 treats failures and selective outcomes as empirical evidence about the architecture rather than as cases to be tuned away after freeze.

\section{Contributions}
\begin{itemize}
    \item \textbf{A literature-grounded D01--D10 operational framework} spanning everyday practices, pragmatics, etiquette, values, institutions, religion, family and gender, work and civic life, heritage, and identity/intergroup relations.
    \item \textbf{A standalone evidence-grounded verifier architecture} that separates prompt context, cultural applicability, response assessability, cultural-dimension planning, deterministic response-span grounding, bounded evidence acquisition, frozen evidence comparison and final cultural scoring.
    \item \textbf{Explicit selective semantics} distinguishing \texttt{not\_culturally\_applicable}, \texttt{not\_assessable} and \texttt{insufficient\_evidence}, while retaining dimension-level \texttt{abstain} instead of fabricating 0/1/2 judgments.
    \item \textbf{A narrow contextual-reasoning fallback} for evidence-unresolved context-dependent recommendations only; unresolved external facts and descriptive norms remain eligible for abstention.
    \item \textbf{A reproducible evidence layer} with source typing, leakage controls, LIVE acquisition, REPLAY over frozen snapshots, stable identifiers and trace validation.
    \item \textbf{A controlled Best-of-4 evaluation mode} used for comparison with five-annotator human preferences, Skywork Reward V2 and a direct LLM judge, without making ranking the primary product interface.
    \item \textbf{A staged empirical program} including the six-source 120-item positive-control run, the frozen PLT comparison and a post-freeze external-validation protocol.
\end{itemize}

\section{Thesis Structure}
This thesis is organized into six chapters. Chapter~2 reviews the literature on cultural appropriateness in large language models, cultural evaluation and adaptation, model-based evaluation, and evidence-grounded verification, and concludes by identifying the research gap addressed in this work. Chapter~3 presents the research design, the operational framework for cultural appropriateness, the Vericult methodology and implementation, and the experimental design. Chapter~4 reports the evaluation results, including comparisons with reward-model and LLM-judge baselines, agreement with human judgments, and external challenge evaluation. Chapter~5 discusses the findings, analyzes errors and implications, and examines limitations and threats to validity. Finally, Chapter~6 summarizes the work and outlines directions for future research.
